\documentclass[12pt,a4paper]{article}
\usepackage[margin=1in]{geometry}
\usepackage{amsmath}
\usepackage{amssymb}
\usepackage{graphicx}
\usepackage{hyperref}
\usepackage{caption}
\usepackage{float}

\title{Lattice QCD Analysis Methods}
\author{HMC GPU Simulation}
\date{\today}

\begin{document}

\maketitle

\section*{1. Autocorrelation Function}

For a sequence of measurements $\{O_0, O_1, O_2, \dots, O_{N-1}\}$, the autocorrelation function is defined as:

\begin{equation}
\Gamma(t) = \frac{1}{N-t} \sum_{i=0}^{N-t-1} (O_i - \langle O \rangle)(O_{i+t} - \langle O \rangle)
\end{equation}

where the average is

\begin{equation}
\langle O \rangle = \frac{1}{N} \sum_{i=0}^{N-1} O_i
\end{equation}

\section*{2. Normalized Autocorrelation Function}

The normalized autocorrelation function is:

\begin{equation}
\rho(t) = \frac{\Gamma(t)}{\Gamma(0)}
\end{equation}

where $\Gamma(0)$ is the variance of the data:

\begin{equation}
\Gamma(0) = \text{Var}(O) = \frac{1}{N} \sum_{i=0}^{N-1} (O_i - \langle O \rangle)^2
\end{equation}

\section*{3. Integrated Autocorrelation Time}

The integrated autocorrelation time characterizes the number of effectively independent measurements:

\begin{equation}
\tau_{\text{int}} = \frac{1}{2} + \sum_{t=1}^{W} \rho(t)
\end{equation}

where $W$ is the cutoff window. In our analysis, we use $W = N/2$.

\section*{4. Statistical Error}

The error on the mean, accounting for autocorrelation, is:

\begin{equation}
\sigma_{\langle O \rangle} = \sqrt{\frac{2\tau_{\text{int}}}{N}} \sigma_O
\end{equation}

where $\sigma_O = \sqrt{\Gamma(0)}$ is the standard deviation.

\section*{5. Plaquette Observable}

The plaquette is the fundamental Wilson loop around a single plaquette:

\begin{equation}
P_{\mu\nu}(x) = \text{Tr}\left[ U_{\mu}(x) U_{\nu}(x+\hat{\mu}) U_{\mu}^{\dagger}(x+\hat{\nu}) U_{\nu}^{\dagger}(x) \right]
\end{equation}

The average plaquette:

\begin{equation}
\langle P \rangle = \frac{1}{N_t N_s^3 \cdot 6\cdot 3} \sum_{x} \sum_{\mu<\nu} \text{Re}\left[ P_{\mu\nu}(x) \right]
\end{equation}

\section*{6. Topological Charge}

The topological charge is computed using the clover field strength tensor:

\begin{equation}
Q = \frac{1}{32\pi^2} \sum_x \sum_{\mu\nu\rho\sigma} 
\varepsilon_{\mu\nu\rho\sigma} \text{Tr}\left[ F_{\mu\nu}(x) F_{\rho\sigma}(x) \right]
\end{equation}

The field strength tensor is constructed from 4 oriented plaquettes at each site:

\begin{equation}
C_{\mu\nu}(x) = \frac{1}{4} \left[ P_{\mu\nu}(x) + P_{\nu,-\mu}(x) + P_{-\mu,-\nu}(x) + P_{-\nu,\mu}(x) \right]
\end{equation}


\begin{equation}
F_{\mu\nu}(x) = \frac{1}{2}(C_{\mu\nu}-C_{\mu\nu}^\dagger)_{traceless} 
\end{equation}

\section*{7.  Smearing}

To improve the resolution of topological structures, we apply Stout smearing iteratively. For each smearing step, the links are updated as follows:

\begin{equation}
U_\mu(x) \to U'_\mu(x) = \frac{1}{\mathcal{N}} \left[ U_\mu(x) + \alpha \, \Sigma_\mu(x) \right]
\end{equation}

where $\Sigma_\mu(x)$ is the staple sum and $\mathcal{N}$ is a normalization factor ensuring $U'_\mu$ remains unitary after reunitarization via Gram-Schmidt.

The staple for direction $\mu$ is constructed from neighboring links:

\begin{equation}
\Sigma_\mu(x) = \sum_{\nu \neq \mu} \left[ U_\nu(x) U_\mu(x+\hat{\nu}) U_\nu^\dagger(x+\hat \mu) + 
U_\nu^\dagger(x-\hat{\nu}) U_\mu(x-\hat{\nu}) U_\nu(x-\hat{\nu}+\hat\mu) \right]
\end{equation}

After computing the updated link, we perform reunitarization to project back to SU(3).

\begin{table}[H]
\centering
\caption{Smearing parameters used in this analysis}
\begin{tabular}{|c|c|c|}
\hline
Parameter & Value & Description \\
\hline
$\alpha$ & 0.1 & Smearing strength \\
$n_{\text{iter}}$ & 8 & Number of smearing steps \\
\hline
\end{tabular}
\end{table}

\section*{8. Simulation Parameters}

\begin{table}[H]
\centering
\caption{Lattice parameters and simulation details}
\begin{tabular}{|c|c|c|c|c|}
\hline
Label & $N_T \times N_S^3$ & $\beta$ & Trajectories & Measurements \\
\hline
T8S4 & $8 \times 4^3$ & 5.7, 6.0, 6.3 & 10000 & 80000 (8 smears $\times$ 10000)\\
T10S8 & $10 \times 8^3$ & 5.7, 6.0, 6.3 & 10000 & 80000 \\
T16S10 & $16 \times 10^3$ & 5.7, 6.0, 6.3 & 10000 & 80000 \\
\hline
\end{tabular}
\end{table}

Each trajectory consists of 10 molecular dynamics steps with $\tau = 1.0$.

\section*{9. Basic Observables Summary}

\begin{table}[H]
\centering
\caption{Summary of plaquette and topological charge measurements at smear=0}
\begin{tabular}{|c|c|c|c|c|c|}
\hline
Lattice & $\beta$ & $\langle P \rangle$ & $\sigma_P$ & $\langle Q \rangle$ & $\sigma_Q$ \\
\hline
$8 \times 4^3$ & 5.7 & 0.5574 & 0.0075 & -0.0003 & 0.1376 \\
$8 \times 4^3$ & 6.0 & 0.5962 & 0.0059 & 0.0023 & 0.1336 \\
$8 \times 4^3$ & 6.3 & 0.6236 & 0.0051 & -0.0009 & 0.1292 \\
$10 \times 8^3$ & 5.7 & 0.5492 & 0.0028 & -0.0072 & 0.4389 \\
$10 \times 8^3$ & 6.0 & 0.5940 & 0.0019 & 0.0071 & 0.4278 \\
$10 \times 8^3$ & 6.3 & 0.6226 & 0.0017 & -0.0002 & 0.4069 \\
$16 \times 10^3$ & 5.7 & 0.5504 & 0.0020 & -0.0012 & 0.7859 \\
$16 \times 10^3$ & 6.0 & 0.5938 & 0.0011 & -0.0147 & 0.7395 \\
$16 \times 10^3$ & 6.3 & 0.6226 & 0.0009 & 0.0048 & 0.7271 \\
\hline
\end{tabular}
\end{table}

\section*{9. Results: Plaquette vs Smearing}

\begin{figure}[H]
\centering
\includegraphics[width=0.8\textwidth]{../data_analyze/basic/plaq_vs_smear.png}
\caption{Average plaquette as a function of smearing steps for different lattice sizes and $\beta$ values.}
\end{figure}

\section*{10. Results: Topological Charge vs Smearing}

\begin{figure}[H]
\centering
\includegraphics[width=0.8\textwidth]{../data_analyze/basic/topo_vs_smear.png}
\caption{Topological charge evolution with smearing steps for different lattice sizes and $\beta$ values.}
\end{figure}

\section*{11. Results: Autocorrelation Analysis}

The autocorrelation function $\rho(t)$ measures the correlation between measurements separated by $t$ MD time units.

\begin{figure}[H]
\centering
\includegraphics[width=0.8\textwidth]{../data_analyze/result1/plaq_autocorr_T8_S4.png}
\caption{Plaquette autocorrelation function for $8 \times 4^3$ lattice at different $\beta$ values.}
\end{figure}

\begin{figure}[H]
\centering
\includegraphics[width=0.8\textwidth]{../data_analyze/result1/plaq_autocorr_T10_S8.png}
\caption{Plaquette autocorrelation function for $10 \times 8^3$ lattice at different $\beta$ values.}
\end{figure}

\begin{figure}[H]
\centering
\includegraphics[width=0.8\textwidth]{../data_analyze/result1/plaq_autocorr_T16_S10.png}
\caption{Plaquette autocorrelation function for $16 \times 10^3$ lattice at different $\beta$ values.}
\end{figure}

\section*{12. Integrated Autocorrelation Times}

\begin{table}[H]
\centering
\caption{Integrated autocorrelation time $\tau_{\text{int}}$ for plaquette measurements at $\beta=6.0$ ($T=8, S=4$)}
\begin{tabular}{|c|c|c|c|}
\hline
Smear & $\beta=5.7$ & $\beta=6.0$ & $\beta=6.3$ \\
\hline
0 & 20.19 & 3.01 & 2.07 \\
1 & 22.22 & 3.25 & 2.31 \\
2 & 24.10 & 3.46 & 2.55 \\
3 & 25.85 & 3.65 & 2.80 \\
4 & 27.50 & 3.82 & 3.06 \\
5 & 29.07 & 3.98 & 3.33 \\
6 & 30.56 & 4.13 & 3.62 \\
7 & 31.99 & 4.28 & 3.91 \\
\hline
\end{tabular}
\end{table}

\begin{table}[H]
\centering
\caption{Integrated autocorrelation time $\tau_{\text{int}}$ for plaquette measurements ($T=10, S=8$)}
\begin{tabular}{|c|c|c|c|}
\hline
Smear & $\beta=5.7$ & $\beta=6.0$ & $\beta=6.3$ \\
\hline
0 & 52.71 & 11.54 & 2.92 \\
1 & 56.19 & 13.10 & 3.37 \\
2 & 59.13 & 14.57 & 3.81 \\
3 & 61.63 & 15.97 & 4.26 \\
4 & 63.83 & 17.33 & 4.70 \\
5 & 65.78 & 18.67 & 5.15 \\
6 & 67.55 & 20.01 & 5.60 \\
7 & 69.18 & 21.39 & 6.07 \\
\hline
\end{tabular}
\end{table}

\begin{table}[H]
\centering
\caption{Integrated autocorrelation time $\tau_{\text{int}}$ for plaquette measurements ($T=16, S=10$)}
\begin{tabular}{|c|c|c|c|}
\hline
Smear & $\beta=5.7$ & $\beta=6.0$ & $\beta=6.3$ \\
\hline
0 & 157.10 & 45.79 & 10.39 \\
1 & 163.48 & 51.50 & 11.12 \\
2 & 168.71 & 57.06 & 11.84 \\
3 & 173.08 & 62.51 & 12.55 \\
4 & 176.80 & 67.91 & 13.27 \\
5 & 180.03 & 73.29 & 14.01 \\
6 & 182.88 & 78.70 & 14.78 \\
7 & 185.42 & 84.13 & 15.58 \\
\hline
\end{tabular}
\end{table}

\section*{13. Topological Charge Autocorrelation}

\begin{table}[H]
\centering
\caption{Integrated autocorrelation time $\tau_{\text{int}}$ for topological charge ($T=8, S=4$)}
\begin{tabular}{|c|c|c|c|}
\hline
Smear & $\beta=5.7$ & $\beta=6.0$ & $\beta=6.3$ \\
\hline
0 & 0.49 & 0.80 & 0.37 \\
1 & 0.50 & 0.86 & 0.37 \\
2 & 0.52 & 0.92 & 0.37 \\
3 & 0.56 & 0.96 & 0.39 \\
4 & 0.61 & 0.99 & 0.43 \\
5 & 0.69 & 1.01 & 0.47 \\
6 & 0.78 & 1.02 & 0.51 \\
7 & 0.89 & 1.01 & 0.56 \\
\hline
\end{tabular}
\end{table}

\begin{table}[H]
\centering
\caption{Integrated autocorrelation time $\tau_{\text{int}}$ for topological charge ($T=10, S=8$)}
\begin{tabular}{|c|c|c|c|}
\hline
Smear & $\beta=5.7$ & $\beta=6.0$ & $\beta=6.3$ \\
\hline
0 & 0.61 & 0.63 & 0.71 \\
1 & 0.66 & 0.66 & 0.78 \\
2 & 0.73 & 0.69 & 0.85 \\
3 & 0.83 & 0.71 & 0.93 \\
4 & 0.95 & 0.74 & 1.01 \\
5 & 1.11 & 0.78 & 1.09 \\
6 & 1.30 & 0.84 & 1.16 \\
7 & 1.52 & 0.92 & 1.25 \\
\hline
\end{tabular}
\end{table}

\begin{table}[H]
\centering
caption{Integrated autocorrelation time $\tau_{\text{int}}$ for topological charge ($T=16, S=10$)}
\begin{tabular}{|c|c|c|c|}
\hline
Smear & $\beta=5.7$ & $\beta=6.0$ & $\beta=6.3$ \\
\hline
0 & 6.76 & 2.81 & 1.53 \\
1 & 7.88 & 3.05 & 1.64 \\
2 & 8.97 & 3.34 & 1.78 \\
3 & 9.99 & 3.65 & 1.93 \\
4 & 10.89 & 3.99 & 2.09 \\
5 & 11.66 & 4.32 & 2.26 \\
6 & 12.30 & 4.64 & 2.43 \\
7 & 12.80 & 4.94 & 2.62 \\
\hline
\end{tabular}
\end{table}

\section*{14. Discussion}

The results show a striking observation: $\tau_{\text{int}}(Q) < \tau_{\text{int}}(P)$ for all parameter sets. This is opposite to the expected QCD behavior where topological charge typically has a longer autocorrelation time than the plaquette.

Possible explanations:
\begin{itemize}
\item The topological charge is frozen in the trivial sector ($Q \approx 0$)
\item The lattice volumes are too small to observe topological transitions
\item The Wilson flow scale $w_0$ may not be properly set
\end{itemize}

At $\beta=6.3$, the plaquette autocorrelation time decreases significantly, indicating the system is in the deconfined phase.

Near the phase transition, the integrated autocorrelation time diverges:

\begin{equation}
\tau_{\text{int}} \propto \xi^z \sim |\beta - \beta_c|^{-z}
\end{equation}

where $\xi$ is the correlation length and $z \approx 2$ for local algorithms.

\end{document}
